\section{Upper Bound on $\|H\|_{r \to r}$ via Riesz--Thorin Interpolation}

We use the following standard result.

\begin{theorem}[Riesz--Thorin interpolation theorem]
Let $(\Omega,\mu)$ be a $\sigma$-finite measure space, and let $1\le p_0,p_1,q_0,q_1\le \infty$. For $0<\theta<1$, define $p_\theta,q_\theta$ by
\[
\frac1{p_\theta} = \frac{1-\theta}{p_0} + \frac{\theta}{p_1}, \qquad \frac1{q_\theta} = \frac{1-\theta}{q_0} + \frac{\theta}{q_1},
\]
with the convention $1/\infty=0$. Suppose $T:L^{p_0}(\Omega)+L^{p_1}(\Omega)\to L^{q_0}(\Omega)+L^{q_1}(\Omega)$ is linear, and that for $i=0,1$,
\[
T\bigl(L^{p_i}(\Omega)\bigr)\subseteq L^{q_i}(\Omega), \qquad \|Tf\|_{q_i}\le M_i\|f\|_{p_i} \quad\text{for all }f\in L^{p_i}(\Omega).
\]
Then $T\bigl(L^{p_\theta}(\Omega)\bigr)\subseteq L^{q_\theta}(\Omega)$ and $\|T\|_{L^{p_\theta}\to L^{q_\theta}} \le M_0^{1-\theta}M_1^\theta$.
\end{theorem}

Let $\Omega=\{1,\dots,m\}$ with counting measure $\mu$. Then $\mu(\Omega)=m<\infty$, so $(\Omega,\mu)$ is $\sigma$-finite, and $L^p(\Omega,\mu)=\ell_p^m$ for every $1\le p\le \infty$.

Define $T:\mathbb{R}^m\to \mathbb{R}^m$ by $Tx=Hx$. Since $L^2(\Omega,\mu;\mathbb{R})+L^\infty(\Omega,\mu;\mathbb{R})=\mathbb{R}^m$ (all $L^p$ spaces coincide on a finite set), $T$ is one linear map on the common domain, and its endpoint actions on $L^2$ and $L^\infty$ are restrictions of the same operator, verifying the compatibility hypothesis. (We note that the Riesz--Thorin theorem applies equally to real-valued function spaces; indeed, the real induced $\ell^p$ operator norm satisfies $\|A\|_{\ell^p(\mathbb{R}^m)\to\ell^p(\mathbb{R}^m)}\leq \|A\|_{\ell^p(\mathbb{C}^m)\to\ell^p(\mathbb{C}^m)}$, so the complex interpolation bound is at least as strong as what we need. In our case, both the upper and lower bounds are achieved by real vectors, so the real and complex norms coincide.)

From Section~1, $\|H\|_{2\to 2}=\sqrt{m}$ and $\|H\|_{\infty\to\infty}=m$. We apply Riesz--Thorin with
\[
p_0=q_0=2, \quad p_1=q_1=\infty, \quad M_0=\sqrt{m}, \quad M_1=m.
\]

For $r=2$, $\|H\|_{2\to 2}=\sqrt{m}=m^{1-1/2}$, and the desired estimate holds.

For $r>2$, choose $\theta=1-2/r$. Since $r>2$, we have $0<\theta<1$. Then
\[
\frac{1}{p_\theta} = \frac{1-\theta}{2} + \frac{\theta}{\infty} = \frac{1-\theta}{2} = \frac{1}{r},
\]
and similarly $1/q_\theta = 1/r$. So $p_\theta = q_\theta = r$. Applying Riesz--Thorin:
\[
\|H\|_{r\to r} \le (\sqrt{m})^{1-\theta} m^\theta = (\sqrt{m})^{2/r} m^{1-2/r} = m^{1/r} m^{1-2/r} = m^{1-1/r}.
\]

Therefore, $\|H\|_{r\to r}\le m^{1-1/r}$ for every real $r\ge 2$.
